\section{Signed Constraints and Fixed Witnesses}
\label{sec:variance}

We remove variance by giving each variable two signed copies. The copies
must denote dual trees; the assignment involution expresses this condition.
We then express each side of a failed query by a \emph{fixed witness}: a
solution of the signed system that is fixed by the assignment involution and
exhibits the failure along a given word. The four-spine construction of
\cref{sec:spines} will enforce such witnesses.

\subsection{The Signed System}\label{sec:signed}
Write $V^{\pm}=V\times\{0,1\}$, $v^+=v^0$, $v^-=v^1$, and define the
\emph{variable flip} $\flip(v^q)=v^{1-q}$. Translate $\phi$ to a covariant constraint system
$\psi=\phi^{\pm}$ as follows:

\begin{center}
\begin{tabular}{@{}lll@{}}
\toprule
Original literal & \multicolumn{2}{c}{Signed literals}\\
\midrule
$z\le f(a_i)_i$ & $z^+\le f(a_i^{c(i)})_i$ & $f(a_i^{1\oplus c(i)})_i\le z^-$\\
$f(a_i)_i\le z$ & $f(a_i^{c(i)})_i\le z^+$ & $z^-\le f(a_i^{1\oplus c(i)})_i$\\
$z=\bot$ & $z^+=\bot$ & $z^-=\top$\\
$z=\top$ & $z^+=\top$ & $z^-=\bot$\\
\bottomrule
\end{tabular}
\end{center}

\begin{lemma}[Signed correspondence]\label{lem:signed}
For an assignment $\sigma$ on signed variables define
$(S\sigma)(v)=\mathfrak d(\sigma(\flip v))$. The map $S$, called the
\emph{assignment involution}, is an involution and preserves solutions of $\psi$.
Moreover,

\begin{equation}\label{eq:g14}
\sigma\longmapsto \sigma^{\pm},\qquad \sigma^{\pm}(v^q)=H_q(\sigma v)
\end{equation}

is a bijection between solutions of $\phi$ and $S$-fixed solutions of
$\psi$.
\end{lemma}

\begin{proof} Tree duality $\mathfrak d$ reverses inequalities and commutes with
$f$. Applying it and $\flip$ exchanges the two literals in each row of
the table. This proves preservation by $S$; applying both involutions twice gives
$S^2=\mathrm{id}$. By \eqref{eq:g9}--\eqref{eq:g10},
\eqref{eq:g14} satisfies each translated pair and is fixed by $S$. Conversely, for
an $S$-fixed solution define $\sigma_0(v)=H_0(\sigma(v^+))$. The fixed-point
identity gives $\sigma(v^-)=\mathfrak d(\sigma(v^+))$, so \eqref{eq:g9} reconstructs \eqref{eq:g14}. Applying \eqref{eq:g10}
to the positive literal of each pair proves $\sigma_0\models_c\phi$.
The reconstruction also proves both inverse identities. \end{proof}


\subsection{Fixed Witnesses}\label{sec:fixed-witness}
Write $\theta\in\{\mathsf l,\mathsf r\}$ for the left and right sides;
$\mathsf l$-unsafe means left unsafe, and similarly on the right. Both sides
read the same signed system $\psi$, from different query roots:
\[
 (X_{\mathsf l},Y_{\mathsf l})=(x^+,y^+),\qquad
 (X_{\mathsf r},Y_{\mathsf r})=(y^-,x^-).
\]
The right side exchanges the roots and uses their negative copies, because
the two copies of a variable denote dual trees: a bottom prefix of $y$ is a
top prefix of $y^-$.

We call a covariant solution $\sigma\models\psi$ a \emph{fixed
witness} for $w$ on side $\theta$ if $S\sigma=\sigma$ and
\begin{equation}\label{eq:fixed-witness}
 \widehat{\sigma(X_\theta)}(w)=\top,\qquad
 \widehat{\sigma(Y_\theta)}(w)\ne\top.
\end{equation}
Here \emph{fixed} means fixed by the assignment involution $S$, not that
the word $w$ is held fixed.

\begin{lemma}[Fixed-witness criterion]\label{lem:fixed-witness}
For either side $\theta$, a word $w$ is $\theta$-unsafe iff it has a fixed
witness on side $\theta$.
\end{lemma}
\begin{proof}
By \cref{lem:signed}, the $S$-fixed solutions of $\psi$ are exactly the
assignments $\sigma_0^{\pm}$ for solutions $\sigma_0$ of $\phi$, with
$\sigma_0^{\pm}(v^+)=H_0(\sigma_0v)$ and, by \eqref{eq:g9},
$\sigma_0^{\pm}(v^-)=H_1(\sigma_0v)=\mathfrak d(H_0(\sigma_0v))$. For the left
side, \eqref{eq:fixed-witness} for $\sigma_0^{\pm}$ therefore says that
$\sigma_0(x)$ has a top prefix along $w$ and $\sigma_0(y)$ has none, which is
the definition of a left-unsafe word for $\sigma_0$. For the right side, tree
duality exchanges the labels $\bot$ and $\top$ at every node, so
$\widehat{\mathfrak d(t)}(w)=\top$ iff $\widehat t(w)=\bot$. Hence
\eqref{eq:fixed-witness} at the roots $y^-$ and $x^-$ says that $\sigma_0(y)$
has a bottom prefix along $w$ and $\sigma_0(x)$ has none, which is the
definition of a right-unsafe word for $\sigma_0$. Both equivalences include
paths that terminate before the end of $w$.
\end{proof}

\begin{example}[The signed running witness]\label{ex:running-signed}
For \cref{ex:running-arrow}, the six signed literals can be grouped as
\[
 \begin{gathered}
 x^+\le f(a^-,b^+),\qquad f(a^-,b^+)\le x^+,
 \qquad y^+\le f(a^-,b^+),\\
 f(a^+,b^-)\le x^-,\qquad x^-\le f(a^+,b^-),
 \qquad f(a^+,b^-)\le y^-.
 \end{gathered}
\]
The contravariant child uses the opposite sign; passing to the negative
root also reverses the inequality. The assignment of
\cref{ex:running-arrow} translates to
\[
\begin{array}{c|cccc}
 v&x&y&a&b\\\hline
 \sigma^{\pm}(v^+)&f(\Top,\Top)&f(\Bot,\Bot)&\Bot&\Top\\
 \sigma^{\pm}(v^-)&f(\Bot,\Bot)&f(\Top,\Top)&\Top&\Bot
\end{array}
\]
and each column consists of dual trees, so $S\sigma^{\pm}=\sigma^{\pm}$.
At $w=0$, the labels at $x^+$ and $y^+$ are $\Top$ and $\Bot$, giving the
left fixed witness. The labels at $y^-$ and $x^-$ are also $\Top$ and
$\Bot$, so the same assignment is the right fixed witness for this word.
\end{example}

A covariant solution need not itself be fixed. \Cref{sec:spines} constructs
a fixed solution from the least shape (\cref{sec:least}) of any consistent
system that is flip-closed and sign coherent (\cref{sec:closure}). Flip
closure means closure under reversing inequalities, flipping variables, and
exchanging constants. Sign coherence means that each constructor child has
the sign of its root, toggled by the variance of its position. Applying that
construction to an extension encoding \eqref{eq:fixed-witness} will preserve
the witness conditions as well as satisfaction.
